\begin{document}
\maketitle

\begin{abstract}
This note documents a Lean~4 formalisation of the algebraic
carrier rigidity theorem of TFPT Paper~2 (\emph{Carrier Rigidity,
Hypercharge, and the Standard-Model Packet from Boundary
Polarization}), up to and including the determinant-normalised
$3+2$ hypercharge carrier
$Y = \operatorname{diag}(-\tfrac13,-\tfrac13,-\tfrac13,\tfrac12,\tfrac12)$.
The bundled top-level theorem consumes four explicit interfaces:
a carrier involution, Higgs and Yukawa rank witnesses, and a
primitive oriented determinant-preserving carrier condition.
Each interface is refinable upstream via a separately formalised
Lean theorem. In v10 each upstream interface is closed by a
further structural layer: the Calder\'on idempotent
$\pi^2 = \pi$ is itself produced from a
\texttt{BoundaryPolarization} ($H = H_+ \oplus H_-$); the Higgs
algebraic shadow is identified with the degree-1 piece of the
bivariate polynomial ring
$K[X_0, X_1]_1 = \texttt{MvPolynomial.homogeneousSubmodule (Fin 2) K 1}$
(the algebraic-geometry "Proj-side" presentation of
$H^0(\mathbb P^1, \mathcal O(1))$); and the Stage D Yukawa
trilinear form is constructed explicitly from any basis of a
rank-$3$ space, with injective and surjective contraction proved
by direct basis computation. The Yukawa Stage D characterisation
$\exists\,\omega\;(\textrm{CI}\,\wedge\,\textrm{CS})
  \iff \dim E = 3$ now holds under \texttt{Nontrivial}\,$E$. The unique integer
pair $(q_-,q_+)=(-2,3)$ follows from a general $(m,n)$
lattice-rigidity theorem applied to the structurally bundled
\texttt{PrimitiveOrientedDeterminantCarrier 3 2} input. Lean
proves the unique $3+2$ hypercharge carrier, trace zero, and the
carrier polynomial $6 Y^2 - Y - \mathbf 1 = \mathbf 0$. The
project builds against Mathlib~v4.29.1 with zero \texttt{sorry},
zero \texttt{admit}, no domain-specific axioms, no
\texttt{unsafe}/\texttt{opaque}/\texttt{partial} escape hatches,
and every load-bearing theorem depends only on the three standard
Lean axioms (\texttt{propext}, \texttt{Classical.choice},
\texttt{Quot.sound}). A seven-check \texttt{audit.sh} gate is
shipped; the headline-theorem checks combine Lean
\texttt{\#check} elaboration (\texttt{AuditCheck.lean}) with
strict \texttt{example : <exact type> := <theorem>} signature
locks (\texttt{AuditContract.lean}), so neither silent renames
nor silent weakening of theorem statements can mask regressions.
\end{abstract}

\begin{tfptfrontbox}
\textbf{Inputs from previous papers.} Paper~1 (boundary primitive
kernel) for the \emph{conceptual} derivation of the carrier
involution from the Calder\'on polarisation. The Lean project
consumes only algebraic data: an involution element $\varepsilon$ with
$\varepsilon^2=\mathbf 1$, plus Stage-A certificates for the Higgs and
Yukawa rank discharges.

\textbf{New contribution.} A reproducible Lean~4 project
\texttt{TfptCarrier} (Mathlib~v4.29.1) that elevates four
previously imported inputs to formal Lean theorems:

\begin{itemize}[leftmargin=1.5em, itemsep=0pt]
    \item boundary involution $\Rightarrow$ orthogonal idempotents;
    \item general integer rigidity $\Rightarrow$ unique primitive pair;
    \item structural trace identity $\Rightarrow$ $\tr Y=0$ from
          rank data alone;
    \item determinant character of the carrier torus
          $\Rightarrow$ Diophantine condition $m q_- + n q_+ = 0$
          as a corollary (universal-$\lambda$ form, with a single-
          point helper at $\lambda=2$).
\end{itemize}

Two further inputs ($\dim E_+=2$ and $\dim E_-=3$) are represented
as structured hypotheses, not bare numerical assumptions. They
mark the exact interfaces where the future geometric (Higgs) and
representation-theoretic (Yukawa) proofs must enter.

For the Yukawa case, this note now provides \emph{four}
interface layers, with the upstream stage carrying genuine
algebraic content:

\begin{itemize}[leftmargin=1.5em, itemsep=1pt]
\item \emph{Stage A} (\texttt{YukawaRankCertificate}) --- a
    bare linear equivalence $E_- \simeq K^3$. Used by the
    bundled \texttt{CarrierData.CarrierPremises} as a rank
    witness; refinable from any of the upstream stages.
\item \emph{Stage B} (\texttt{YukawaTopForm}) --- a linear
    equivalence $\Lambda^3 E_- \simeq K$ (an algebraic top
    form on $E_-$).
\item \emph{Stage C} (\texttt{PrimitiveIndecomposableYukawaCoupling})
    --- a linear equivalence $E_- \simeq (\Lambda^2 E_-)^*$
    plus non-triviality (the contraction map of a non-degenerate
    alternating trilinear coupling, taken as data).
\item \emph{Stage D} (\texttt{YukawaTrilinearForm}) --- introduced
    in v6, with operational predicates renamed in v7 --- an
    alternating trilinear form $\omega : \Lambda^3 E_- \to K$ as
    primary datum, together with two explicit \texttt{Prop}
    predicates \texttt{ContractionInjective} and
    \texttt{ContractionSurjective}. From these the contraction
    map is \emph{constructed} (via
    \texttt{alternatingMapLinearEquiv} and
    \texttt{AlternatingMap.curryLeft}) and proved to be a linear
    iso when both predicates hold and $E_-$ is non-trivial,
    producing Stage C as a \emph{theorem}, not as a definition.
\end{itemize}

Each upstream stage produces the next via a Lean bridge
(D $\Rightarrow$ C $\Rightarrow$ B $\Rightarrow$ A). Stage D
carries the genuine content of TFPT Paper 2's
"non-degenerate primitive alternating trilinear coupling": the
trilinear form is the primary datum, the contraction iso is a
derived fact, and the rank-3 conclusion follows from the
contraction iso via the elementary dimension chain
$\dim E_- = \dim (\Lambda^2 E_-)^* = \binom{\dim E_-}{2}$ on a
non-trivial finite-dimensional space.

A seven-check \texttt{audit.sh} CI gate is included.

\textbf{Not claimed here.} No analytic content of Paper~1, no full
geometric formalisation of $H^0(\mathbb P^1, \mathcal O(1))$, no
quiver-theoretic definition of \emph{primitive indecomposable
Yukawa type}, no SM hypercharge spectrum on derived
representations (duals, exterior powers, tensor products), no
$\mathbb Z_6$ quotient, no $\alpha$, no flavour, no cosmology.

\textbf{Falsification or audit surface.} The note fails if
(i)~\texttt{lake build} no longer succeeds against the pinned
Mathlib release; (ii)~the \texttt{audit.sh} script returns a
non-zero exit; (iii)~the dependency graph stated below drifts from
the wording of Paper~2~\S2.
\end{tfptfrontbox}

\begin{tfptclaimcontract}
\textbf{Claim.}
The top-level Lean theorem
\texttt{CarrierPremises.hypercharge\_carrier\_packet} takes
four explicit interfaces:

\begin{enumerate}[leftmargin=1.5em, itemsep=0pt, label=(\arabic*)]
\item a carrier involution element $\varepsilon$ with
      $\varepsilon^2 = \mathbf 1$
      (\texttt{InvolutionProjectors.CarrierInvolution}, optionally
      wrapped by a \texttt{CalderonCertificate});
\item a Higgs rank-2 certificate
      (\texttt{HiggsIndexCertificate}) discharging $\dim E_+ = 2$;
\item a Yukawa rank-3 witness
      (\texttt{YukawaRankCertificate}) discharging $\dim E_- = 3$,
      \emph{refinable} via \texttt{YukawaTopForm} (Stage B),
      \texttt{PrimitiveIndecomposableYukawaCoupling} (Stage C),
      and \texttt{YukawaTrilinearForm} (Stage D) by formal Lean
      bridges;
\item a primitive oriented determinant-preserving carrier
      condition
      $(3 q_- + 2 q_+ = 0,\ q_- < 0 < q_+,\ \gcd(|q_-|,q_+) = 1)$,
      packaged as
      \texttt{PrimitiveOrientedDeterminantCarrier 3 2}.
\end{enumerate}

\textbf{Conclusion.}
From these inputs Lean derives the $3+2$ hypercharge carrier
$Y = -\tfrac13 P_- + \tfrac12 P_+$, with $(q_-,q_+) = (-2, 3)$,
$\tr Y = 0$, and the carrier polynomial
$6 Y^2 - Y - \mathbf 1 = \mathbf 0$ in $\mathrm{End}_K M$.

\textbf{Three layers of the contract.}
The claim is structured into three explicit layers, separated to
prevent overclaiming and to make every open question visible:

\smallskip
\noindent
\emph{(a) Direct inputs of \texttt{CarrierPremises}.}\ The
bundled top-level theorem consumes the four interfaces above:
a carrier involution, a Higgs rank witness, a Yukawa rank
witness, and a primitive oriented determinant-preserving
carrier condition.

\smallskip
\noindent
\emph{(b) Lean-proved upstream refinements (closed to the
next algebraic layer).}\ Each of the four direct inputs is
produced by an explicit Lean-proved upstream interface:

\begin{itemize}[leftmargin=1.6em, itemsep=0pt]
  \item Carrier involution $\leftarrow$
    \texttt{CalderonProjector} $\leftarrow$
    \texttt{BoundaryPolarization} ($H = H_+ \oplus H_-$).
    The converse \texttt{BoundaryPolarization.ofCarrierInvolution}
    makes the chain a Lean-proved \emph{algebraic
    equivalence}: the carrier involution and the polarisation
    are two presentations of the same finite-dimensional
    carrier datum. This is not the analytic Paper-1 content;
    the analytic construction of the polarisation from
    boundary elliptic theory remains upstream.
  \item Higgs rank witness $\leftarrow$
    \texttt{HiggsTopForm} (Stage B) $\leftarrow$
    \texttt{HiggsSchemeCohomologyShadow}
    ($\simeq \texttt{MvPolynomial.homogeneousSubmodule (Fin 2) K 1}$).
  \item Yukawa rank witness $\leftarrow$
    \texttt{YukawaTopForm} (Stage B) $\leftarrow$
    \texttt{YukawaPrimitive} (Stage C) $\leftarrow$
    \texttt{YukawaTrilinearForm} (Stage D) $\leftrightarrow$
    \texttt{YukawaStageDExistence} (\emph{both directions
    formalised, under} \texttt{Nontrivial}).
  \item Primitive oriented determinant carrier $\leftarrow$
    \texttt{SeamWindingInterface} (algebraic packaging of
    seam-class data).
\end{itemize}

\smallskip
\noindent
\emph{(c) Open TFPT-origin theorems (not formalised here).}\ Three
genuinely upstream questions remain open:

\begin{enumerate}[leftmargin=1.6em, itemsep=0pt, label=(\arabic*)]
  \item \emph{Paper 1 analytic boundary theory} produces the
    polarisation $H = H_+ \oplus H_-$ (boundary elliptic
    theory, Sobolev spaces, Calder\'on operator).
  \item \emph{TFPT boundary kernel} canonically produces an
    $\omega : \Lambda^3 E_- \to K$ with
    \texttt{ContractionInjective} and
    \texttt{ContractionSurjective} (the \emph{canonical}
    question; the existence is now closed by the present
    note).
  \item \emph{Seam winding $[u_\Sigma] = 1$} produces a
    primitive oriented determinant-preserving carrier
    condition.
\end{enumerate}

These three are \emph{outside the scope} of the present Lean
note; each is named, typed, and provided with a Lean-side
target structure
(\texttt{BoundaryPolarization},
\texttt{BoundaryYukawaKernelInterface},
\texttt{SeamWindingInterface}).

\textbf{Not claimed.}
Beyond the three open TFPT-origin theorems, the
$\mathbb Z_6$ quotient $G_{\mathrm{phys}}$, derived SM
representations (duals, exterior powers, tensor products),
$\alpha$, flavour, and cosmology are also not formalised
in this note.

\textbf{First assumptions on $K$.}
The projector construction requires $2$ invertible in the
carrier algebra; for vector-space applications this is ensured
by $\mathrm{char}\,K \ne 2$. The rational normalisation
$Y = -P_-/3 + P_+/2$ further requires $6$ invertible; concretely
the bundled theorems use \texttt{CharZero K}.

\textbf{Proof status.}
\emph{Formally verified} in Lean~4 + Mathlib~v4.29.1; no
\texttt{sorry}; no \texttt{admit}; no
\texttt{unsafe}/\texttt{opaque}/\texttt{partial}; no
domain-specific axiom; \texttt{audit.sh} (seven checks
including the signature-lock \texttt{AuditContract.lean})
returns \texttt{AUDIT:\ PASS}. As of v10, the three upstream
interfaces (Calder\'on, Higgs, Yukawa Stage D backward) are
each closed \emph{to the next algebraic layer}. The remaining
gaps are TFPT-side or analytic, not Lean-side.

\textbf{Kill condition.}
Any of the seven audit checks fail; any headline theorem in
\texttt{AuditCheck.lean} is removed or fails to elaborate;
any \texttt{example} in \texttt{AuditContract.lean} fails to
typecheck (which detects silent weakening of theorem statements).
\end{tfptclaimcontract}

\tableofcontents

% ---------------------------------------------------------------- %
\section{What This Note Is About}
% ---------------------------------------------------------------- %

v1 formalised the \emph{algebraic motor} of TFPT Paper~2 ---
the carrier polynomial and the trace identity on a concrete
$5\times5$ diagonal model. v2 began making the upstream
dependencies explicit. v3 made the dependency graph correct.
v4 replaced the Yukawa rank witness with a structurally
algebraic Stage B (top form on $\Lambda^3 E_-$). v5 introduced
a Stage C contraction iso $E_- \simeq (\Lambda^2 E_-)^*$.

v6 completed the Yukawa refinement chain by introducing a
Stage D primary datum; v7 renamed its predicates to operational
names; v8 added a Calder\'on idempotent layer and a Higgs Stage B;
v9 consolidates the document into a single review-ready state
and adds an unconditional transparency theorem for Stage D.
The Stage D primary datum reads:

\[
\texttt{YukawaTrilinearForm}\,K\,E_-
\;\;:=\;\;
\bigl\{\,\omega : \Lambda^3 E_- \to_{\ell K} K\,\bigr\},
\qquad
\texttt{ContractionInjective}\,\omega,\;\;
\texttt{ContractionSurjective}\,\omega.
\]

The contraction map $C_\omega : E_- \to (\Lambda^2 E_-)^*$ is
now \emph{constructed} from $\omega$ via
\texttt{exteriorPower.alternatingMapLinearEquiv} composed with
\texttt{AlternatingMap.curryLeft}. The two predicates state
exactly that $C_\omega$ is injective (mathematical reading:
non-degeneracy) and surjective (mathematical reading: primitivity
/ no spectator). Together, on a non-trivial finite-dimensional
space, they yield a Stage C structure as a theorem; the rank-3
conclusion then follows mechanically:

\[
\boxed{\;
\omega\;\text{trilinear} \;+\; \texttt{CI} \;+\; \texttt{CS} \;+\; \texttt{Nontrivial}\,E_-
\;\;\xRightarrow{\text{contraction is iso}}\;\;
\text{Stage C}
\;\;\xRightarrow{\binom{n}{2} = n,\, n\ne 0}\;\;
\dim E_- = 3.
\;}
\]

\noindent
The non-triviality hypothesis is essential: for $E_- = 0$ both
predicates are vacuously satisfied (the contraction $0 \to 0$ is
trivially bijective) but $\dim E_- = 0$. The Lean theorem
\texttt{finrank\_E\_eq\_zero\_or\_three} (new in v9) records the
unconditional disjunction $\dim E_- = 0 \lor \dim E_- = 3$ for
transparency; the rank-3 slogan above is the case under
\texttt{Nontrivial}\,$E_-$.

The full Yukawa tower in Lean is now
\[
\text{Stage D}
\;\Rightarrow\; \text{Stage C}
\;\Rightarrow\; \text{Stage B}
\;\Rightarrow\; \text{Stage A}
\;\Rightarrow\; \dim E_- = 3,
\]
with every step a Lean theorem and Stages B and C carrying
characterisation theorems establishing equivalence to
$\dim E_- = 3$.

v6 made the Paper~1 input explicit. A new
\texttt{CalderonInterface.lean} introduces a
\texttt{CalderonCertificate} structure that names the
Calder\'on-induced carrier involution as a Lean datum and feeds
it into the \texttt{InvolutionProjectors} chain. v8 upgrades
this interface from documentary to a structural algebraic step:
a new \texttt{CalderonProjector.lean} module shows that an
idempotent element $\pi$ in a ring with $2$ invertible yields
$\varepsilon^2 = 1$ via the algebraic identity $\varepsilon := 2\pi - 1$,
$(2\pi - 1)^2 = 4\pi^2 - 4\pi + 1 = 1$. v10 adds the next
upstream layer: a new \texttt{BoundaryPolarization.lean} module
shows that any direct-sum decomposition $H = H_+ \oplus H_-$
of a $K$-module produces an idempotent endomorphism on $H$ (via
\texttt{Submodule.IsCompl.projection}) and feeds the Calder\'on
chain. The remaining Paper-1 gap is therefore precisely the
analytic content: producing the boundary spectral
decomposition itself.

v8 also introduces a Stage-B Higgs interface
(\texttt{HiggsTopForm.lean}) parallel to the Yukawa tower:
the rank-2 conclusion $\dim E_+ = 2$ is derivable from a
structural hypothesis identifying $E_+$ with the algebraic
shadow of $H^0(\mathbb P^1, \mathcal O(1))$, namely the space
of homogeneous degree-1 polynomials in two formal variables
$\alpha X + \beta Y$. v10 closes the bridge to actual algebraic
geometry: \texttt{HiggsSchemeCohomologyShadow.lean} provides a
\texttt{LinearEquiv} between this algebraic shadow and the
Mathlib-native degree-1 piece of the bivariate polynomial ring
$\texttt{MvPolynomial.homogeneousSubmodule (Fin 2) K 1}
= K[X_0, X_1]_1$, which by the standard $\operatorname{Proj}$
construction is the algebraic-geometry presentation of
$H^0(\mathbb P^1_K, \mathcal O(1))$. The remaining gap is now
the scheme-theoretic identification
$H^0(\mathrm{Proj}\,K[X, Y], \mathcal O(1)) \simeq K[X, Y]_1$,
a general algebraic-geometry theorem independent of TFPT.

v10 also closes the Stage D backward direction
(\texttt{YukawaStageDExistence.lean}). Given any basis
$b : \texttt{Basis}\,(\texttt{Fin}\,3)\,K\,E$, the determinant
form $b.\texttt{det}$ transports via
\texttt{exteriorPower.alternatingMapLinearEquiv} to an
alternating trilinear form $\omega : \Lambda^3 E \to_{\ell K} K$.
The contraction $C_\omega$ is shown injective by basis-level
matrix-determinant computation and surjective by
$\dim E = \dim(\Lambda^2 E)^* = 3$ and the standard
finite-dimensional injective-iff-surjective theorem. From this
the Stage D characterisation
$\exists\,\omega\;(\textrm{CI}\,\wedge\,\textrm{CS})
  \iff \dim E = 3$ now holds (under $\texttt{Nontrivial}\,E$).

The bundled main theorem packages the three load-bearing
conclusions
\[
(q_-,q_+)=(-2,3), \qquad \tr Y=0, \qquad 6\,Y^2 - Y - \mathbf 1 = \mathbf 0
\]
into a single Lean statement \texttt{hypercharge\_carrier\_packet}
that consumes only structured hypotheses --- with the Yukawa
input now refinable all the way to Stage D, and the boundary
input now arriving via an explicit Calder\'on certificate.

The redesign is governed by one editorial discipline: \emph{the
language must not outrun the formal status}. Two examples:

\begin{itemize}[leftmargin=1.6em, itemsep=2pt]
\item This note never claims to formalise the \emph{Standard-
    Model hypercharge spectrum}, only the $3+2$ carrier
    multiplet. Duals, exterior powers, tensor products --- the
    constructions that lift the carrier to all SM chiral
    multiplets --- are out of scope (\S\ref{sec:limits}).
\item The Yukawa Stage-A certificate is named
    \texttt{YukawaRankCertificate}, not
    \texttt{PrimitiveYukawaTypeCertificate}. The certificate is
    a \emph{rank witness}; it is not a formalisation of
    "primitive indecomposable Yukawa type", which is the proper
    open problem.
\end{itemize}

% ---------------------------------------------------------------- %
